本文针对通用神经网络处理器上依赖图切分、同核任务合并与共享只读缓存联动的多核调度问题，建立从合法方案构造到官方事件评估的统一模型，依据操作—张量依赖、私有存储容量、跨核同步和共享带宽约束，采用有限候选划分与逐例成对比较，得到题定配置下最高的五核平均加速比\textbf{3.6104（无量纲）}。研究以全局完成周期最小化为目标。三个场景均覆盖全部官方算例和一至五核配置，结果同时保留分位分布、候选差异与时间线核验，以避免仅凭局部负载或单一搬运量解释全局性能。

\textbf{针对问题一}，场景A规定每个子图是独立Task，私有缓存不能跨Task沿用，因而切分过细会增加DDR中转与固定等待。建立操作覆盖、收缩依赖无环及同核顺序相容的调度模型，以官方全局Makespan为目标函数，约束核心编号、子图划分和任务先后。采用\textbf{连通块分核}，并与依赖深度分段、宽图分批两个候选比较；每个候选先过合法性校验，再使用未改动的官方评估器择优。得到五核平均加速比\textbf{3.5093（无量纲）}，其分位范围与结构分层共同表明\textbf{多数图可利用多个核心}，但大块占优或新增传输较高的图保留低加速尾部。独立重算后，所选五百组方案的覆盖、时间线及官方响应交叉核验全部通过；逐例结果仍显示增核不必然单调加速。

\textbf{针对问题二}，基于问题一的合法切图，在场景B中把同核子图合并为Task，使中间张量可留在L1或UB，同时承担容量压力与可能的Spill。建立场景B调度模型，沿用全局周期最小化目标，并引入逐时刻私有存储容量、跨核边界COPY及固定同步约束，采用\textbf{场景B重新择优}，并把问题一选中方案直接放入B作同方案A/B对照，以分开硬件规则和方案选择的收益。得到五核平均加速比\textbf{3.5456（无量纲）}；成对时长比的逐例范围揭示\textbf{Task合并并不保证逐例加速}，其原因可能是驻留量、Spill与共享DDR争用共同改变关键路径。独立重算后，所选五百组方案的覆盖、时间线及官方响应交叉核验全部通过，容量峰值和新增搬运分解按官方详细输出逐例对照。

\textbf{针对问题三}，进一步在场景B中引入按逻辑张量编号查询的共享只读L2；命中与未命中分别竞争Cache读池和DDR读池，FIFO次序与容量影响后续复用。建立两读池、缓存状态与依赖事件共同决定Makespan的调度模型，以全局周期最小化为目标，并约束只在输入搬运事件使用L2。采用\textbf{同方案反事实}比较题定缓存开关，得到五核平均加速比\textbf{3.6104（无量纲）}及同方案缓存时长比均值\textbf{1.0260（无量纲）}；\textbf{有正向收益的算例过半}，但命中字节率不能直接等同于完成时间收益。所选五百组方案的覆盖、时间线及官方响应交叉核验全部通过；另在六个有代表性的案例上改变容量和带宽，独立验证默认格与正式结果一致，并将该网格明确作为机制敏感性分析。

创新点在于先保持切图更换Task语义、再保持方案单独切换L2读路径；相比只报告各场景最终方案的速度，成对设计另以\textbf{1.0260（无量纲）}的五核同方案Cache比度量缓存机制效果。

